\clearpage
\section{Lamb Shanks in Massaman Curry}

%\index{author!Nagi Maehashi}
\indexbyrecipe{source!recipetineats.com}
\indexbyrecipe{meat!lamb}
\indexbyrecipe{starch!potato}
\indexbyrecipe{starch!rice}
%\index{country!Thailand}
%\index{curry!Thai}
%\index{stew}

Source: https://www.recipetineats.com/lamb-shanks-in-massaman-curry/

Serves: 5

\subsection{Ingredients}

\begin{itemize}
    \item 1.5kg lamb shanks (5 small, 4 medium or 2--3 large)
    \item 114g Massaman curry paste
    \item 400mL full-fat coconut milk
    \item 2 cups low-sodium chicken stock
    \item 1 onion, halved and sliced 1cm thick
    \item 400g small potatoes, halved if necessary
    \item 1 star anise
    \item 1 cinnamon stick
    \item Sliced red chilli and coriander, to garnish
    \item Steamed jasmine rice, to serve
\end{itemize}

\subsection{Method}

\begin{enumerate}
    \item Preheat the oven to \celsius{180} (\celsius{160} fan).
    \item Mix the curry paste, coconut milk and stock in a $22 \times 33$cm baking dish. Add the onion, potatoes, star anise, cinnamon and lamb shanks.
    \item Turn the shanks through the sauce until coated, then cover the dish tightly with foil.
    \item Bake for 2 hours. Remove the foil and continue baking for 1 hour if using small shanks, or 1\sfrac{1}{2} hours for medium or large shanks. Turn the shanks twice during this time so they brown evenly.
    \item When the meat is tender enough to pull apart easily, transfer the shanks to a plate. Skim excess fat from the sauce; it should have reduced to a slightly syrupy consistency.
    \item Serve the shanks and sauce with jasmine rice, garnished with chilli and coriander.
\end{enumerate}

\subsection{Notes}

Use low-sodium stock because the sauce reduces considerably. If the sauce becomes too thick, loosen it with a little boiling water. If it remains thin after the lamb is tender, return the sauce alone to the oven to reduce.
